\documentclass[twocolumn,a4paper]{article}

\usepackage[T1]{fontenc}
\usepackage[a4paper, top=2.5cm, bottom=2.5cm, left=2cm, right=2cm]{geometry}
\usepackage{newtxtext,newtxmath}
\usepackage{amsmath}
\usepackage{graphicx}
\usepackage{booktabs}
\usepackage{multirow}
% \usepackage{svg}  % requires Inkscape; SVGs are pre-converted to PDF
\usepackage{parskip}
\usepackage{float}
\usepackage{xcolor}
\usepackage{titlesec}
\usepackage{abstract}
\usepackage{hyperref}

\usepackage[style=apa, sortcites=true, sorting=ynt]{biblatex}
\addbibresource{referencias.bib}

\graphicspath{{../RECURSOS/figures/}{../RECURSOS/images/}}

% ── Title spacing ───────────────────────────────────────────────
\titlespacing*{\section}{0pt}{1.2ex plus .5ex minus .2ex}{0.8ex}
\titlespacing*{\subsection}{0pt}{1ex plus .4ex minus .2ex}{0.5ex}

% ── Abstract spans both columns ─────────────────────────────────
\renewcommand{\abstractnamefont}{\normalfont\bfseries}
\renewcommand{\abstracttextfont}{\normalfont\small}

\begin{document}

% ── Title block (single column) ──────────────────────────────────
\twocolumn[{%
  \begin{@twocolumnfalse}
    \title{\large\bfseries Draft: Improvements on Community Membership Hiding\\
    via a Custom Reinforcement Learning Architecture}
    \author{}
    \date{}
    \maketitle
    \begin{abstract}
      \noindent
      We address the problem of Community Membership Hiding (CMH): given a graph
      and a community detection algorithm, strategically rewire a target node's
      edges so it is no longer grouped into its original community.
      Building on the Deep Reinforcement Learning (DRL) baseline of
      \parencite{bernini-2024}, we identify two fundamental architectural
      limitations—target-unawareness and static node embeddings—and propose
      successive improvements culminating in a \textbf{Dual-Agent GATv2
      Actor-Critic} that specializes in edge addition and deletion separately.
      We further introduce a \textit{Policy-Guided Subgraphing} heuristic that
      makes the architecture tractable on large graphs. Empirical evaluation on
      five benchmark datasets shows consistent and substantial gains in F1 score
      over all baselines, including the gradient-based methods of
      \parencite{silvestri2025righthidemaskingcommunity}.
    \end{abstract}
    \vspace{1em}
  \end{@twocolumnfalse}%
}]


% ════════════════════════════════════════════════════════════════
\section{Introduction}
% ════════════════════════════════════════════════════════════════

The modern science of networks has brought significant advances to our
understanding of complex systems. One of the most relevant features of graphs
representing real-world structures—ranging from social media platforms to
biological networks—is their community structure: tightly connected groups of
nodes identified by highly effective algorithms
\parencite{Fortunato-2010}. However, in some contexts this functionality
inherently risks exposing individuals to privacy breaches by inadvertently
revealing sensitive affiliations, such as political leanings or religious
beliefs, based solely on their structural position.

To counter this, the concept of Community Membership Hiding (CMH) has been
proposed \parencite{bernini-2024}. The objective of CMH is to strategically
alter the structural properties of a network so that a specific target node is
no longer grouped into its original community by a detection algorithm, while
maintaining overall graph integrity and treating the algorithm as a black box.

Recent work has formulated this evasion task as a Markov Decision Process (MDP)
solved via Deep Reinforcement Learning (DRL) \parencite{bernini-2024}. While
that work demonstrated the viability of the approach, it exhibits notable
performance degradation on larger graphs and has fundamental architectural
limitations that hinder optimal performance. This paper develops directly from
\parencite{bernini-2024}, dissects its shortcomings, and proposes a novel RL
architecture that leverages modern Graph Neural Networks (GNNs). We compare
against three strong baselines: the original agent of \parencite{bernini-2024},
and the gradient-based methods $\nabla$-CMH and $\nabla$-CMH (projected) from
\parencite{silvestri2025righthidemaskingcommunity}.


% ════════════════════════════════════════════════════════════════
\section{Community Evasion}
% ════════════════════════════════════════════════════════════════

\subsection{Problem Definition}

The following definitions are borrowed from \parencite{bernini-2024}. Let
$\mathcal{G} = (V, E)$ be a graph where $V$ is the set of nodes and $E$ is the
set of edges. Given a community detection algorithm $f(\cdot)$ and a target
node $u \in V$, $f$ assigns $u$ to an original community $\mathcal{C}_i$.

The goal of CMH is to find a modified graph $\mathcal{G}' = (V, E')$ such that
upon re-evaluating $f(\mathcal{G}')$ the new community assigned to $u$, denoted
$\mathcal{C}_i'$, is sufficiently dissimilar to the original. Formally, $u$ is
considered hidden when:
\begin{equation}
  \operatorname{sim}(\mathcal{C}_i - \{u\},\; \mathcal{C}_i' - \{u\}) \leq \tau
  \label{eq:hiding}
\end{equation}
where $\tau \in [0,1]$ is a predefined hiding threshold, and
$\operatorname{sim}(\cdot,\cdot)$ is a similarity function (we use the
Sørensen-Dice coefficient, following the baseline).

To aid hiding, the action space is restricted: the agent may only remove edges
from $u$ to nodes $v \in \mathcal{C}_i$, and add edges from $u$ to nodes
$v \notin \mathcal{C}_i$. The maximum number of modifications (edge budget) is:
\[
  B = \max\!\left(1,\left\lfloor \tfrac{|E|}{|V|} \cdot \beta \right\rfloor\right)
\]

Beyond hiding success, a desired property is to maximize Normalized Mutual
Information (NMI)—a measure of similarity between the original and modified
community structures—thereby minimizing structural distortion. We report the
\textbf{F1 score}, defined as the harmonic mean of the Success Rate (SR) and
NMI, as the main evaluation metric.

\subsection{Original Agent Architecture}

The baseline of \parencite{bernini-2024} uses a Graph Convolutional Network
(GCN) \parencite{kipf2017semisupervisedclassificationgraphconvolutional} to
process graph state. Initial node features are generated with \textit{node2vec}
\parencite{grover2016node2vecscalablefeaturelearning} on the unaltered graph.
The agent is constrained to act on the target node via an external environment
wrapper that masks invalid actions.

\subsection{Architectural Limitations of the Baseline}

Two critical bottlenecks limit the baseline's scalability and effectiveness:

\textbf{Target-unawareness.} The GCN receives the entire graph as input but
lacks explicit knowledge of the target node $u$. Because standard GNNs are
permutation-invariant, the policy $\pi(a \mid s)$ produces identical outputs
for any target node given the same graph state. The agent therefore learns a
generic ``shuffle communities'' strategy rather than a node-specific one.

\textbf{Static embeddings.} \textit{Node2vec} embeddings are computed once on
$\mathcal{G}$ and never updated. As the agent rewires the graph, these
embeddings diverge from the current topology, creating an information lag that
degrades learning efficiency.


% ════════════════════════════════════════════════════════════════
\section{Proposed Architectures}
% ════════════════════════════════════════════════════════════════

\subsection{GATv2 Actor-Critic v1}
\textcolor{blue}{TODO: agree on a name for this version.}

Our first proposed architecture introduces an \textbf{Actor-Critic} model with
a \textbf{shared graph encoder} based on Graph Attention Networks v2 (GATv2)
\parencite{brody2022attentivegraphattentionnetworks}. Instead of static
node2vec embeddings, we define a \textit{target-aware} feature space that is
partially recomputed after each rewiring step.

\textbf{Static features} (computed once on the original graph):
\begin{itemize}
  \item \textit{Is Target} — identifies the node being hidden.
  \item \textit{In Target Community} — marks the original deception boundary.
  \item \textit{Is Original Neighbor} — preserves the initial neighborhood context.
\end{itemize}

\textbf{Dynamic features} (updated after every rewiring action):
\begin{itemize}
  \item \textit{Normalized degree} — tracks local density changes.
  \item \textit{Current neighbor status} — binary edge indicator w.r.t.\ the target.
  \item \textit{Common neighbors} — structural overlap with the target.
  \item \textit{Jaccard similarity} to the target.
\end{itemize}

At the core, a \textbf{shared GNN backbone} of stacked GATv2 layers encodes
node features and edge attributes into node embeddings $H \in \mathbb{R}^{N
\times h}$. These embeddings are consumed by both heads:

\begin{itemize}
  \item \textbf{Actor head.} Computes per-node logits conditioned on the
    target embedding $h_t$ via a scoring function
    $\ell_i = f(W_1 h_i + W_2 h_t)$, where $h_i$ is the candidate node
    embedding. The target node's own logit is masked to $-\infty$.

  \item \textbf{Critic head.} Combines an attention-pooled global graph
    representation $h_G = \sum_i \alpha_i h_i$ (with
    $\alpha_i = \operatorname{softmax}(g(h_i))$) with the target embedding
    to produce a scalar state-value estimate $V(\mathcal{G}, u)$.
\end{itemize}

% ── Figure placeholder ──────────────────────────────────────────
\begin{figure}[h]
  \centering
  \fbox{\rule{0pt}{4.5cm}\rule{\linewidth}{0pt}}
  \caption{\textit{[Placeholder]} GATv2 Actor-Critic v1 architecture diagram.}
  \label{fig:ac-v1}
\end{figure}

The reward signal follows the CMH objective: a positive terminal reward on
successful hiding, a negative terminal reward on budget exhaustion, and
step-level costs that differ for edge addition vs.\ deletion to encourage
efficient rewiring.


\subsection{GATv2 Actor-Critic v2 (Dual-Agent)}
\textcolor{blue}{TODO: agree on a name for this version.}

To further improve performance—particularly NMI preservation and scalability—we
move from a single monolithic agent to a \textbf{Dual-Agent Architecture} that
decomposes the rewiring task into two specialized sub-problems.

\subsubsection*{Agent Specialization}

\begin{itemize}
  \item \textbf{Adder Agent.} Trained exclusively to identify and create optimal
    inter-community edges from $u$ to nodes outside $\mathcal{C}_i$.
  \item \textbf{Deleter Agent.} Trained exclusively to identify and remove
    critical intra-community edges from $u$ to nodes inside $\mathcal{C}_i$.
\end{itemize}

Both agents share the same GATv2-based backbone design (same feature space,
same encoder architecture), but maintain \emph{independent} actor and critic
heads, each trained on their respective action subspace.

\subsubsection*{Action-Selection Mechanism}

At each timestep, both agents are queried. The agent selected to act is the one
whose Critic reports the higher expected state value:
\[
  \text{act} = \arg\max_{a \in \{\text{add}, \text{del}\}} V_a(\mathcal{G}, u)
\]
This value-based selection allows the architecture to adaptively switch between
addition and deletion depending on which is more beneficial given the current
graph state, without requiring a hand-crafted switching rule.

% ── Figure placeholder ──────────────────────────────────────────
\begin{figure}[h]
  \centering
  \fbox{\rule{0pt}{5cm}\rule{\linewidth}{0pt}}
  \caption{\textit{[Placeholder]} Dual-Agent GATv2 Actor-Critic architecture.
    Both agents share a backbone design; action selection is governed by
    comparing Critic values.}
  \label{fig:ac-v2}
\end{figure}

\subsubsection*{Scaling via Policy-Guided Subgraphing}

A key bottleneck on large graphs is the cost of recomputing dynamic features
(Common Neighbors, Jaccard Similarity) for all nodes after every action, which
grows with $|V|$ and $|E|$. We address this with a \textbf{Policy-Guided
Subgraphing} heuristic:

\begin{enumerate}
  \item \textbf{Initial pass.} Perform a single forward pass of both Actor
    networks on the full graph to obtain raw logits.
  \item \textbf{Top-node identification.} Select the top-$k$ nodes most
    preferred by each agent (by logit magnitude).
  \item \textbf{Community expansion.} Retrieve the full communities of each
    selected top node.
  \item \textbf{Subgraph construction.} Build a reduced subgraph containing
    only those communities and the target node.
  \item \textbf{Local operation.} Dynamic feature updates and all subsequent
    actions operate entirely within this subgraph.
\end{enumerate}

This heuristic trades a small, controllable reduction in F1 score for a
dramatic decrease in per-step computation, making previously intractable large
graphs feasible. We compare it against a simpler \textit{k-neighborhood}
baseline (extracting the $k$-hop ego-graph of $u$), which misses distant
``bridge'' nodes and performs worse in practice.

% ── Figure placeholder ──────────────────────────────────────────
\begin{figure}[h]
  \centering
  \fbox{\rule{0pt}{3.5cm}\rule{\linewidth}{0pt}}
  \caption{\textit{[Placeholder]} Policy-Guided Subgraphing: full graph
    (left) reduced to a task-relevant subgraph (right) guided by agent logits.}
  \label{fig:subgraph}
\end{figure}


% ════════════════════════════════════════════════════════════════
\section{Overlapping Communities}
% ════════════════════════════════════════════════════════════════

\textcolor{blue}{[This section is left intentionally blank — to be filled in.]}


% ════════════════════════════════════════════════════════════════
\section{Results}
\label{results}
% ════════════════════════════════════════════════════════════════

We report results using the F1 score (harmonic mean of SR and NMI), following
the standard CMH evaluation protocol \parencite{bernini-2024}. Experiments were
run on an AWS EC2 \texttt{g4dn.xlarge} instance (4 Intel Xeon vCPUs, 16 GiB
RAM, NVIDIA T4 GPU) over five benchmark datasets of varying size and density.
We swept $\tau \in \{0.3, 0.5, 0.8\}$ and $\beta \in \{0.5, 1.0, 2.0\}$.

\subsection{GATv2 Actor-Critic v1 vs.\ Baselines}

Table~\ref{tab:success_rate_comparison} compares success rates of GATv2 v1
against the original agent across all $(\tau, \beta)$ combinations. Our
architecture maintains a consistent and substantial lead across all
configurations and does not overfit to any single scenario.

\begin{table}[h]
\centering
\small
\begin{tabular}{llcc}
\toprule
\textbf{$\beta$} & \textbf{$\tau$} & \textbf{GATv2} & \textbf{Original} \\
\midrule
\multirow{3}{*}{0.5} & 0.3 & 61.67 & 41.67 \\
                     & 0.5 & 70.21 & 38.22 \\
                     & 0.8 & 79.38 & 62.67 \\
\midrule
\multirow{3}{*}{1.0} & 0.3 & 79.17 & 61.67 \\
                     & 0.5 & 91.78 & 58.44 \\
                     & 0.8 & 96.25 & 90.00 \\
\midrule
\multirow{3}{*}{2.0} & 0.3 & 82.08 & 77.67 \\
                     & 0.5 & 98.75 & 74.67 \\
                     & 0.8 & 99.17 & 96.33 \\
\bottomrule
\end{tabular}
\caption{Success rates (\%) for GATv2 v1 vs.\ original agent
         ($\tau=0.5,\ \beta=1$ shown in Figure~\ref{fig:tradeoff-f1-time}).}
\label{tab:success_rate_comparison}
\end{table}

Figure~\ref{fig:tradeoff-f1-time} shows the F1/time tradeoff for
$\tau=0.5,\ \beta=1$ across all datasets. GATv2 v1 consistently outperforms
all baselines in F1 score, achieving up to 40\% improvement over $\nabla$-CMH
on the largest graph. The gradient-based methods are faster but sacrifice F1.

% The SVG file tradeoff_plot_compact.svg should be in the same folder.
\begin{figure}[h]
  \centering
  \includegraphics[width=\linewidth]{tradeoff_plot}
  \caption{F1 score vs.\ execution time ($\tau=0.5,\ \beta=1$).
           Time axis in logarithmic scale.}
  \label{fig:tradeoff-f1-time}
\end{figure}

\subsection{Dual-Agent Results}

\textcolor{blue}{[Results for the Dual-Agent (GATv2 v2) to be filled in.]}

% ── Figure placeholder ──────────────────────────────────────────
\begin{figure}[h]
  \centering
  \fbox{\rule{0pt}{4.5cm}\rule{\linewidth}{0pt}}
  \caption{\textit{[Placeholder]} Dual-Agent F1 results across datasets and
    community detection algorithms ($\tau=0.5,\ \beta=1$).}
  \label{fig:dual-results}
\end{figure}

% ── Figure placeholder ──────────────────────────────────────────
\begin{figure}[h]
  \centering
  \fbox{\rule{0pt}{4cm}\rule{\linewidth}{0pt}}
  \caption{\textit{[Placeholder]} Effect of Policy-Guided Subgraphing on
    execution time and F1 score on large graphs.}
  \label{fig:subgraph-scaling}
\end{figure}


% ════════════════════════════════════════════════════════════════
\printbibliography
% ════════════════════════════════════════════════════════════════

\end{document}
